\chapter{De \OO-lijn: E-101, V-103, K-102 en de RWZI}
\label{h:o2}
De zuurstof verlaat de stack \emph{nat}: verzadigd of bijna verzadigd met waterdamp, bij \VAR{n(I.T_O2_uit,0)}\degC\
en \VAR{n(I.P_O2_uit,1)}\,barg. Het gas gaat naar de RWZI. Eerst moet het water eruit, want waterdamp is ballast
voor de compressor, en condensaat in leidingen en compressor geeft corrosie, waterslag en schade. Alles in dit
hoofdstuk geldt \textbf{per unit}.

\section{Stap 1: samenstelling van stroom S13}
\kop{Fysisch}
S13 is een mengsel van \OO\ en waterdamp. Voor condensatie telt niet de totale druk, maar de
\emph{partiële} druk van de damp (Dalton). Die volgt uit de molfractie.
\kop{Gegeven}
$\md_\mathrm{mix} = \VAR{n(I.m_O2mix,0)}\kgh$, $w_{\HHO} = \VAR{n(I.w_H2O,2)}$, $P = \VAR{n(I.P_O2_uit,1)}$\,barg,
$T = \VAR{n(I.T_O2_uit,0)}\degC$. Molaire massa's: $M_{\OO} = \VAR{n(R.M_O2,3)}$ en $M_{\HHO} = \VAR{n(R.M_H2O,3)}$\,kg/kmol
\VAR{"(CoolProp; Cengel A-1 geeft 31,999 en 18,015)" if A else "(Cengel tabel A-1)"}.
\kop{Invullen}
\begin{align*}
\md_{\HHO} &= w\cdot\md_\mathrm{mix} = \VAR{n(I.w_H2O,2)}\cdot\VAR{n(I.m_O2mix,0)} = \VAR{n(R.m_H2O_in,2)}\kgh, &
\md_{\OO} &= \VAR{n(I.m_O2mix,0)} - \VAR{n(R.m_H2O_in,2)} = \VAR{n(R.m_O2,2)}\kgh,\\
N_{\OO} &= \frac{\VAR{n(R.m_O2,2)}}{\VAR{n(R.M_O2,3)}} = \VAR{n(R.N_O2,4)}\,\mathrm{kmol/h}, &
N_{\HHO} &= \frac{\VAR{n(R.m_H2O_in,2)}}{\VAR{n(R.M_H2O,3)}} = \VAR{n(R.N_H2O,4)}\,\mathrm{kmol/h},\\
y_v &= \frac{N_{\HHO}}{N_{\OO}+N_{\HHO}} = \frac{\VAR{n(R.N_H2O,4)}}{\VAR{n(R.N_tot,4)}} = \VAR{n(R.y_v,4)}, &
P &= 0{,}1\cdot100+101{,}325 = \VAR{n(R.P_O2_abs,3)}\kPa,\\
P_v &= y_v\,P = \VAR{n(R.y_v,4)}\cdot\VAR{n(R.P_O2_abs,3)} = \VAR{n(R.P_v,3)}\kPa .
\end{align*}
\begin{waarom}[molfractie en geen massafractie]
De wet van Dalton werkt met \emph{aantallen moleculen}: elk molecuul draagt evenveel bij aan de druk. Water is
lichter dan \OO\ ($18$ tegen $32$). Daarom is 15\,\% water in massa ruim \VAR{n(R.y_v*100,0)}\,\% in mol.
\end{waarom}

\section{Stap 2: dauwpunt van S13}
\begin{methode}[– Cengel H14, dauwpunttemperatuur]
$T_\mathrm{dauw} = T_\mathrm{sat}(P_v)$: de temperatuur waarbij de damp bij zijn eigen partiële druk begint te
condenseren als je het gas afkoelt bij constante druk.
\end{methode}
%% if A
\begin{coolprop}
\texttt{PropsSI("T","P",\VAR{n(R.P_v*1000,0)},"Q",1,"Water")} $-\,273{,}15$ $=$ \textbf{\VAR{n(R.T_dauw,2)}\degC}.
Twee eigenschappen: $P$ en $x=1$ (verzadigde damp). In Cengel is dat tabel A-5 (druktabel) bij
$P = \VAR{n(R.P_v,2)}\kPa$. Met de temperatuurtabel A-4, lineair tussen 65\degC\ (25,043\kPa) en 70\degC\ (31,202\kPa):
$T = 65 + \frac{\VAR{n(R.P_v,3)}-25{,}043}{31{,}202-25{,}043}\cdot 5 = \VAR{n(65 + (R.P_v-25.043)/(31.202-25.043)*5,2)}\degC$.
Verschil: \VAR{n(R.T_dauw - (65 + (R.P_v-25.043)/(31.202-25.043)*5),2)}\,K. De oorzaak is lineaire interpolatie op
een kromme $P_\mathrm{sat}(T)$-lijn, die exponentieel verloopt.
\end{coolprop}
%% else
\begin{tabelkader}[A-4]
$P_v = \VAR{n(R.P_v,3)}\kPa$ ligt tussen $P_\mathrm{sat}(65\degC) = 25{,}043\kPa$ en $P_\mathrm{sat}(70\degC) = 31{,}202\kPa$.
Lineair interpoleren:
\[
T_\mathrm{dauw} = 65 + \frac{\VAR{n(R.P_v,3)} - 25{,}043}{31{,}202 - 25{,}043}\cdot 5 = \VAR{n(R.T_dauw,2)}\degC .
\]
\end{tabelkader}
%% endif
\begin{resultaat}
Het dauwpunt is $\VAR{n(R.T_dauw,1)}\degC$. S13 heeft 70\degC\ en is dus \VAR{n(I.T_O2_uit - R.T_dauw,1)}\,K
\emph{oververhit}: er condenseert nog niets in leiding S13, zolang die geïsoleerd is. De condensatie begint in
E-101, op het moment dat het gas onder $\VAR{n(R.T_dauw,1)}\degC$ komt.
\end{resultaat}

\section{Stap 3: hoeveel damp gaat er mee na de KO-drum V-103?}
\kop{Fysisch}
E-101 koelt het gas tot $T_\mathrm{KO} = \VAR{n(I.T_KO,0)}\degC$. Bij die temperatuur kan het gas maar een beperkte
hoeveelheid damp bevatten. Daarboven condenseert het water. In V-103 scheiden gas en druppels. Het gas dat V-103
verlaat (S17) is \emph{verzadigd} bij 30\degC.
\kop{Methode}
Na V-103 geldt $P_v = P_\mathrm{sat}(30\degC)$, dus $y_{v,\max} = P_\mathrm{sat}(T_\mathrm{KO})/P$. Uit
$y = N_v/(N_v + N_{\OO})$ volgt $N_v = \frac{y}{1-y}\,N_{\OO}$. De \OO\ condenseert niet, dus $N_{\OO}$ blijft gelijk.
\kop{Invullen}
%% if A
\begin{coolprop}
$P_\mathrm{sat}(30\degC)$ $=$ \texttt{PropsSI("P","T",303.15,"Q",0,"Water")} $=$ \VAR{n(R.Psat_KO*1000,1)}\,Pa
$=$ \textbf{\VAR{n(R.Psat_KO,4)}\kPa}. Twee eigenschappen: $T$ en $x=0$. Cengel tabel A-4: $4{,}2469\kPa$.
Het verschil is alleen afronding.
\end{coolprop}
%% else
\begin{tabelkader}[A-4]
$P_\mathrm{sat}(30\degC) = 4{,}2469\kPa$.
\end{tabelkader}
%% endif
\begin{align*}
y_{v,\max} &= \frac{\VAR{n(R.Psat_KO,4)}}{\VAR{n(R.P_O2_abs,3)}} = \VAR{n(R.y_v_max,5)}, \qquad
N_v = \frac{\VAR{n(R.y_v_max,5)}}{1-\VAR{n(R.y_v_max,5)}}\cdot\VAR{n(R.N_O2,4)} = \VAR{n(R.N_v,4)}\,\mathrm{kmol/h},\\
\md_v &= \VAR{n(R.N_v,4)}\cdot\VAR{n(R.M_H2O,3)} = \VAR{n(R.m_v,2)}\kgh \quad\text{(damp mee in S17)},\\
\md_\mathrm{cond} &= \VAR{n(R.m_H2O_in,2)} - \VAR{n(R.m_v,2)} = \VAR{n(R.m_cond,2)}\kgh \quad\text{(condensaat S15)},\\
\md_\mathrm{gas} &= \VAR{n(R.m_O2,2)} + \VAR{n(R.m_v,2)} = \VAR{n(R.m_gas,2)}\kgh \quad\text{(S17)}.
\end{align*}
Uitgedrukt als vochtverhouding: vóór E-101 is $\omega = \VAR{n(R.m_H2O_in,2)}/\VAR{n(R.m_O2,2)} = \VAR{n(R.omega_in,4)}$
en na V-103 is $\omega = \VAR{n(R.m_v,2)}/\VAR{n(R.m_O2,2)} = \VAR{n(R.omega_KO,4)}$ kg water per kg \OO.
\begin{resultaat}
Van de \VAR{n(R.m_H2O_in,1)}\kgh\ water condenseert \VAR{n(R.m_cond,1)}\kgh\
(\VAR{pct(R.m_cond/R.m_H2O_in,0)}). Het condensaat is zuiver (gecondenseerde damp, UPW-kwaliteit) en gaat via
P-101 terug naar de unit.
\end{resultaat}
\begin{waarom}[30\degC\ in de KO-drum]
Hoe kouder, hoe minder damp mee naar K-102: minder compressorarbeid, minder kans op condensatie in de
leiding naar de RWZI en meer water terug naar de unit. De ondergrens wordt gezet door het koelwater
(\VAR{n(I.T_kw_in,0)} $\rightarrow$ \VAR{n(I.T_kw_uit,0)}\degC). Met tegenstroom moet de gasuitlaat (30\degC)
boven de koelwaterinlaat (20\degC) blijven. 10\,K is een gangbaar temperatuurverschil (\emph{approach}).
Kouder dan 30\degC\ vraagt een koelmachine en levert weinig op: van 30 naar 25\degC\ scheelt maar
\VAR{n(WA.basis.m_v - WA.KO25.m_v,1)}\kgh\ damp (bijlage~\ref{b:vragen}).
\end{waarom}

\section{Stap 4: warmtebalans van E-101}
\kop{Fysisch}
E-101 doet twee dingen tegelijk. Het koelt de \OO\ van 70 naar 30\degC\ (voelbare warmte), en het koelt de damp en
laat een deel ervan condenseren (latente warmte). De latente warmte is veruit het grootste deel.
\begin{methode}[– Cengel H5 (warmtewisselaar) + H14 (koelen met ontvochtigen)]
\[
\Qd_{\text{E-101}} = \underbrace{\frac{\md_{\OO}}{3600}\,(h_{\OO,30}-h_{\OO,70})}_{\text{\OO\ voelbaar}}
\;+\; \frac{1}{3600}\Big[\underbrace{\md_v\,h_{g,30}}_{\text{damp die blijft}} + \underbrace{\md_\mathrm{cond}\,h_{f,30}}_{\text{condensaat}}
- \underbrace{\md_{\HHO}\,h_{g,70}}_{\text{damp in}}\Big].
\]
Voor waterdamp bij lage partiële druk (een paar kPa) is de enthalpie die van een ideaal gas. Die hangt alleen van
$T$ af, dus $h_v(T) \approx h_g(T)$ (Cengel H14). Het condensaat verlaat de wisselaar als verzadigde vloeistof
($h_f$).
\end{methode}
%% if A
\begin{coolprop}
{\footnotesize\begin{tabular}{@{}lll@{}}
grootheid & CoolProp-aanroep & waarde\\\hline
$h_{\OO}(70\degC,\,p_{\OO}=\VAR{n(R.pO2_in,2)}\kPa)$ & \texttt{PropsSI("H","T",343.15,"P",\VAR{n(R.pO2_in*1000,0)},"Oxygen")} & \VAR{n(R.hO2_70,3)}\kJkg\\
$h_{\OO}(30\degC,\,p_{\OO}=\VAR{n(R.pO2_ko,2)}\kPa)$ & \texttt{PropsSI("H","T",303.15,"P",\VAR{n(R.pO2_ko*1000,0)},"Oxygen")} & \VAR{n(R.hO2_30,3)}\kJkg\\
$h_v(70\degC)$ & \texttt{PropsSI("H","T",343.15,"P",500,"Water")} & \VAR{n(R.hv_70,2)}\kJkg\\
$h_v(30\degC)$ & \texttt{PropsSI("H","T",303.15,"P",500,"Water")} & \VAR{n(R.hv_30,2)}\kJkg\\
$h_f(30\degC)$ & \texttt{PropsSI("H","T",303.15,"Q",0,"Water")} & \VAR{n(R.hf_30,2)}\kJkg\\
\end{tabular}}

\textbf{Waarom de \OO\ bij partiële druk?} Volgens Dalton ``ziet'' elke component alleen zijn eigen partiële druk.
\textbf{Waarom de damp bij 500\,Pa?} Dat is een kunstgreep om het ideale-gasdeel van de damp te krijgen ($h$ hangt dan
alleen van $T$ af, precies Cengels aanname $h_v \approx h_g(T)$). Cengel A-4 geeft $h_g(70\degC) = 2626{,}1$ en
$h_g(30\degC) = 2555{,}6\kJkg$. CoolProp bij verzadiging geeft \VAR{n(R.hg_70_sat,1)} en \VAR{n(R.hg_30_sat,1)}; bij
500\,Pa \VAR{n(R.hv_70,1)} en \VAR{n(R.hv_30,1)}. Het verschil is kleiner dan 0,5\,kJ/kg.
De effectieve $c_p$ van \OO\ is $(\VAR{n(R.hO2_70,3)}-\VAR{n(R.hO2_30,3)})/40 = \VAR{n(R.cp_O2_eff,4)}\kJkgK$; Cengel A-2 geeft 0,918.
\end{coolprop}
\kop{Invullen}
\begin{align*}
\Qd_{\OO} &= \frac{\VAR{n(R.m_O2,2)}}{3600}\,(\VAR{n(R.hO2_30,3)} - \VAR{n(R.hO2_70,3)}) = \VAR{n(R.Q_c_O2,2)}\kW,\\
\Qd_\mathrm{water} &= \frac{\VAR{n(R.m_v,2)}\cdot\VAR{n(R.hv_30,2)} + \VAR{n(R.m_cond,2)}\cdot\VAR{n(R.hf_30,2)} - \VAR{n(R.m_H2O_in,2)}\cdot\VAR{n(R.hv_70,2)}}{3600} = \VAR{n(R.Q_c_w,2)}\kW,\\
\Qd_{\text{E-101}} &= \VAR{n(R.Q_c_O2,2)} + (\VAR{n(R.Q_c_w,2)}) = \mathbf{\VAR{n(R.Q_E101,2)}\kW}.
\end{align*}
%% else
\begin{tabelkader}[A-2 en A-4]
$c_{p,\OO} = 0{,}918\kJkgK$ (A-2, 300\,K); $h_g(70\degC) = 2626{,}1$, $h_g(30\degC) = 2555{,}6$, $h_f(30\degC) = 125{,}74\kJkg$ (A-4).
\end{tabelkader}
\kop{Invullen}
\begin{align*}
\Qd_{\OO} &= \frac{\VAR{n(R.m_O2,2)}}{3600}\cdot 0{,}918\cdot(30-70) = \VAR{n(R.Q_c_O2,2)}\kW,\\
\Qd_\mathrm{water} &= \frac{\VAR{n(R.m_v,2)}\cdot 2555{,}6 + \VAR{n(R.m_cond,2)}\cdot 125{,}74 - \VAR{n(R.m_H2O_in,2)}\cdot 2626{,}1}{3600} = \VAR{n(R.Q_c_w,2)}\kW,\\
\Qd_{\text{E-101}} &= \VAR{n(R.Q_c_O2,2)} + (\VAR{n(R.Q_c_w,2)}) = \mathbf{\VAR{n(R.Q_E101,2)}\kW}.
\end{align*}
%% endif
\begin{resultaat}
E-101 haalt \VAR{n(-R.Q_E101,1)}\kW\ uit het gas. Daarvan is \VAR{n(R.Q_E101_lat,1)}\kW\
(\VAR{pct(R.Q_E101_lat/(-R.Q_E101),0)}) condensatiewarmte. Het afkoelen van de \OO\ zelf is maar
\VAR{n(-R.Q_c_O2,1)}\kW.
\end{resultaat}
\begin{waarom}[E-101 op koud koelwater en niet op TW]
Het gas moet naar 30\degC. TW komt terug met \VAR{n(R.T_tw_r,0)}\degC\ en is dus te warm: de tweede hoofdwet
laat het gas nooit kouder worden dan het koelmiddel. Daarom koelt E-101 met koelwater van 20\degC. De warmte van
E-101 komt wel indirect terug in de ITM-energiebalans (hoofdstuk~\ref{h:itm}), maar wordt niet teruggewonnen.
\end{waarom}

\kop{Koelwaterdebiet KW1}
%% if A
Met $\Delta h_\mathrm{kw} = h(\VAR{n(I.T_kw_uit,0)}\degC) - h(\VAR{n(I.T_kw_in,0)}\degC) = \VAR{n(R.dh_kw,3)}\kJkg$
(CoolProp, vloeibaar water bij 3\,bar):
%% else
Met $c_{p,w} = 4{,}18\kJkgK$ en $\Delta T = \VAR{n(I.T_kw_uit - I.T_kw_in,0)}$\,K, dus $\Delta h = \VAR{n(R.dh_kw,1)}\kJkg$:
%% endif
\[
\md_\mathrm{KW1} = \frac{-\Qd_{\text{E-101}}\cdot 3600}{\Delta h_\mathrm{kw}} = \frac{\VAR{n(-R.Q_E101,2)}\cdot3600}{\VAR{n(R.dh_kw,3)}} = \VAR{n(R.m_kw_E101,0)}\kgh .
\]

\section{Stap 5: compressie in K-102}
\kop{Fysisch}
Het verzadigde gas (S17) gaat van \VAR{n(R.P_O2_abs,1)}\kPa\ (0,1\,barg) naar \VAR{n(R.P_uit_K102,1)}\kPa\
(\VAR{n(I.P_RWZI,1)}\,barg). Drukverhouding $r_p = \VAR{n(R.rp_K102,3)}$.
\begin{waarom}[0,7\,barg naar de RWZI]
De zuurstof wordt op de RWZI via bodembeluchters in het water geblazen. De druk moet de waterkolom boven de
beluchters overwinnen (5 à 6\,m water $\approx$ 0,5 à 0,6\,bar), plus de verliezen in de beluchterelementen en de
leiding. 0,7\,barg is daarvoor een gangbare ontwerpwaarde (aanname van het model).
\end{waarom}
%% if A
\begin{methode}[– Cengel H7 + H13]
Voor een ideaal mengsel is de mengselentropie $s_\mathrm{mix} = \sum m_i\,s_i(T,P_i)/m$, met elke component bij
zijn partiële druk. Isentroop: zoek $T_{2s}$ zodat $s_\mathrm{mix}(T_{2s},P_2) = s_\mathrm{mix}(T_1,P_1)$. Daarna
$h_2 = h_1 + (h_{2s}-h_1)/\eta$ en $T_2$ uit $h_\mathrm{mix}(T_2) = h_2$.
\end{methode}
\begin{coolprop}
Partiële drukken: in $p_{\OO,1} = \VAR{n(R.pO2_ko,2)}$ en $p_{v,1} = \VAR{n(R.pv_1,3)}\kPa$; uit
$p_{\OO,2} = \VAR{n(R.pO2_2,2)}$ en $p_{v,2} = \VAR{n(R.pv_2,3)}\kPa$ (de molfractie blijft gelijk).
\texttt{PropsSI("S","T",T,"P",p\_O2,"Oxygen")} voor de \OO.
Voor de damp: $s_v(T,p_v) = s(T,500\,\mathrm{Pa}) - R_v\ln(p_v/500)$. Dat is Cengels formule
$\Delta s = -R\ln(P_2/P_1)$ voor een ideaal gas bij constante $T$, met $R_v = R_u/M_{\HHO}$.
De vergelijking $s_\mathrm{mix}(T_{2s}) = s_1$ wordt opgelost met bisectie: steeds het interval halveren
tot $T$ op $10^{-6}$\,K nauwkeurig is.
\end{coolprop}
\kop{Invullen}
\begin{align*}
h_1 &= h_\mathrm{mix}(30\degC) = \VAR{n(R.h1,2)}\kJkg, \qquad s_1 = \VAR{n(R.s1,4)}\kJkgK,\\
T_{2s} &= \VAR{n(R.T2s,2)}\degC \;\Rightarrow\; h_{2s} = \VAR{n(R.h2s,2)}\kJkg,\\
h_2 &= \VAR{n(R.h1,2)} + \frac{\VAR{n(R.h2s,2)} - \VAR{n(R.h1,2)}}{\VAR{n(I.eta_K102,2)}} = \VAR{n(R.h2,2)}\kJkg \;\Rightarrow\; T_2 = \mathbf{\VAR{n(R.T_uit_K102,2)}\degC},\\
\Wd_{\text{K-102}} &= \frac{\VAR{n(R.m_gas,2)}}{3600}\,(\VAR{n(R.h2,2)} - \VAR{n(R.h1,2)}) = \mathbf{\VAR{n(R.W_K102,2)}\kW}.
\end{align*}
\begin{coolprop}[– check met Cengel]
Met ideaal gas en constante $c_p$ (versie B): $T_2 = \VAR{n(RB.T_uit_K102,2)}\degC$ en $\Wd = \VAR{n(RB.W_K102,2)}\kW$. CoolProp geeft
\VAR{n(R.T_uit_K102,2)}\degC\ en \VAR{n(R.W_K102,2)}\kW, een verschil van minder dan 1\,\%. De oorzaak is vooral dat de
handmethode $k_{\OO} = 1{,}395$ gebruikt voor het hele mengsel, terwijl de damp een lagere $k$ heeft ($\approx 1{,}33$).
\end{coolprop}
%% else
\begin{methode}[– Cengel H7, ideaal gas met constante $c_p$]
$T_{2s} = T_1\,r_p^{(k-1)/k}$, $T_2 = T_1 + (T_{2s}-T_1)/\eta$ en $\Wd = \md\,c_{p,\mathrm{mix}}\,(T_2-T_1)$.
\end{methode}
\begin{tabelkader}[A-2]
$c_{p,\OO} = 0{,}918$, $c_{p,\HHO\,\mathrm{damp}} = 1{,}8723\kJkgK$, $k_{\OO} = 1{,}395$ (300\,K).
\end{tabelkader}
\begin{aanname}
Het mengsel is voor \VAR{n(R.m_O2/R.m_gas*100,1)}\,\% (massa) zuurstof. Daarom gebruiken we $k$ van \OO\ voor het hele
mengsel. De damp zou $k$ iets verlagen.
\end{aanname}
\kop{Invullen}
\begin{align*}
c_{p,\mathrm{mix}} &= \frac{\VAR{n(R.m_O2,2)}\cdot 0{,}918 + \VAR{n(R.m_v,2)}\cdot 1{,}8723}{\VAR{n(R.m_gas,2)}} = \VAR{n(R.cp_m,4)}\kJkgK,\\
T_{2s} &= 303{,}15\cdot\VAR{n(R.rp_K102,4)}^{\,0{,}395/1{,}395} = \VAR{n(R.T2s,2)}\,\mathrm{K} = \VAR{n(R.T2s - T0,2)}\degC,\\
T_2 &= 303{,}15 + \frac{\VAR{n(R.T2s,2)} - 303{,}15}{\VAR{n(I.eta_K102,2)}} = \VAR{n(R.T2,2)}\,\mathrm{K} = \mathbf{\VAR{n(R.T_uit_K102,2)}\degC},\\
\Wd_{\text{K-102}} &= \frac{\VAR{n(R.m_gas,2)}}{3600}\cdot\VAR{n(R.cp_m,4)}\cdot(\VAR{n(R.T2,2)} - 303{,}15) = \mathbf{\VAR{n(R.W_K102,2)}\kW}.
\end{align*}
%% endif
\begin{waarom}[$\eta = 0{,}70$]
K-102 is een kleine machine (een blower of schroefcompressor van ruim 10\,kW) met een lage drukverhouding. Bij
kleine machines wegen lekkage en wrijving zwaarder, en 0,65 tot 0,75 is gebruikelijk. Het is een aanname van het
model; met een leveranciersopgave vervang je die.
\end{waarom}

\section{Stap 6: pakketkoeler PK-102 en de stroom naar de RWZI}
De pakketkoeler koelt het gas \VAR{n(I.dT_pakket,0)}\,K af: $T_\mathrm{S18} = \VAR{n(R.T_uit_K102,2)} - \VAR{n(I.dT_pakket,0)} = \VAR{n(R.T_S18,2)}\degC$.
%% if A
\[
\Qd_{\text{PK-102}} = \frac{\md_\mathrm{gas}}{3600}\,\big(h_2 - h_\mathrm{mix}(T_\mathrm{S18})\big) = \VAR{n(R.Q_pak_K102,2)}\kW,\qquad
\Qd_\mathrm{RWZI} = \Wd_{\text{K-102}} - \Qd_{\text{PK-102}} = \VAR{n(R.Q_O2_RWZI,2)}\kW.
\]
%% else
\[
\Qd_{\text{PK-102}} = \md c_p\,\Delta T = \VAR{n(R.mcp_O2,4)}\cdot\VAR{n(I.dT_pakket,0)} = \VAR{n(R.Q_pak_K102,2)}\kW,\qquad
\Qd_\mathrm{RWZI} = \md c_p\,(T_\mathrm{S18}-30) = \VAR{n(R.mcp_O2,4)}\cdot\VAR{n(R.T_S18-30,2)} = \VAR{n(R.Q_O2_RWZI,2)}\kW.
\]
Hierin is $\md c_p = \VAR{n(R.m_gas,2)}/3600\cdot\VAR{n(R.cp_m,4)} = \VAR{n(R.mcp_O2,4)}$\,kW/K.
%% endif
$\Qd_\mathrm{RWZI}$ is de warmte die met het gas meegaat naar de RWZI, gerekend ten opzichte van 30\degC.

\kop{Dauwpunt van S18}
Na compressie stijgt de partiële dampdruk mee met de totale druk:
$p_{v,2} = y_{v,\max}\cdot\VAR{n(R.P_uit_K102,1)} = \VAR{n(R.pv_2,3)}\kPa$.
%% if A
Dan is $T_\mathrm{dauw} = T_\mathrm{sat}(p_{v,2}) = \VAR{n(R.T_dauw_S18,1)}\degC$ (CoolProp, $P$ en $x=1$).
%% else
Interpoleren in A-4 tussen 30\degC\ (4,2469\kPa) en 40\degC\ (7,3851\kPa) geeft $T_\mathrm{dauw} = \VAR{n(R.T_dauw_S18,1)}\degC$.
%% endif
\begin{waarom}[waarom dit ertoe doet]
S18 is $\VAR{n(R.T_S18,1)}\degC$, dus \VAR{n(R.T_S18 - R.T_dauw_S18,0)}\,K boven het dauwpunt. Koelt de leiding naar de
RWZI het gas onder \VAR{n(R.T_dauw_S18,0)}\degC\ af (lange leiding, winter), dan condenseert er water. Daarom is
een aflaat (drain) op het laagste punt nodig, of isolatie. Het dauwpunt na compressie is \emph{hoger} dan in de
KO-drum (30\degC): comprimeren verhoogt $p_v$, terwijl $\omega$ gelijk blijft.
\end{waarom}
\begin{waarom}[PK-102]
Hete zuurstof is gevaarlijker: metalen en afdichtingen ontsteken in \OO\ makkelijker bij een hogere temperatuur.
Lagere temperaturen in \OO-leidingen zijn daarom altijd beter. De warmte (\VAR{n(R.Q_pak_K102,1)}\kW\ per unit) is
klein en wordt niet teruggewonnen.
\end{waarom}
